\section{Trabajos relacionados}
\label{sec:related}

\subsection{Planificación clásica de rutas}
El problema de la ruta de costo mínimo sobre una malla está resuelto desde hace
décadas. El algoritmo de Dijkstra \cite{dijkstra1959} da el óptimo exacto en tiempo
polinomial, y sus variantes con heurística lo aceleran. La condición es fuerte:
hace falta el grafo completo y un mundo determinista. Nosotros no competimos con
Dijkstra en ese terreno; lo usamos como \emph{línea base exacta} para medir cuánto
pierde un agente que solo ve un parche local y que además debe tolerar transiciones
estocásticas, dos supuestos que la planificación clásica no cubre.

\subsection{RL tabular sobre grillas}
El gridworld es el caso de estudio estándar del Aprendizaje por Refuerzo
\cite{sutton2018}, y Q-learning tabular lo resuelve bien cuando el número de estados
es chico. En nuestro trabajo parcial seguimos esa línea: malla $54\times54$, inicio
y meta fijos, una tabla $Q$ con una entrada por par (celda, acción). El límite
aparece al crecer el problema. Con la malla completa de $108\times108$ y un inicio
que cambia en cada episodio, la tabla tendría que almacenar del orden de $10^8$
entradas y, peor aún, no compartiría nada entre estados parecidos: memoriza rutas en
lugar de aprender a navegar. Este trabajo ataca justo ese límite.

\subsection{Métodos profundos basados en valor}
DQN \cite{mnih2015} reemplaza la tabla por una red neuronal y estabiliza el
entrenamiento con dos mecanismos: reproducción de experiencias y una red objetivo
congelada. Double DQN \cite{vanhasselt2016} corrige después el sesgo optimista que
introduce el operador $\max$. Liu y Zou \cite{liu2018} midieron el peso real de la
memoria de reproducción apagándola y variando su tamaño, con juegos de Atari como
banco de pruebas. Nosotros repetimos ese estudio de ablación, pero sobre un modelo
digital de elevación real en lugar de píxeles de videojuego, y con un espacio de
observación de baja dimensión (83 números): un escenario donde no era obvio que los
dos mecanismos siguieran siendo indispensables.

\subsection{Métodos basados en política}
REINFORCE \cite{williams1992} optimiza la política directamente por gradiente, sin
pasar por una función de valor, y es el punto de partida de la familia
actor-critic \cite{mnih2016}. Las comparaciones publicadas entre ambas familias
suelen hacerse sobre entornos de referencia (Atari, MuJoCo). Nuestro aporte aquí es
más modesto y más específico: corremos DQN y REINFORCE sobre el mismo terreno, con
el mismo conjunto de pares de test y la misma métrica, para mostrar el contraste de
estabilidad y varianza en un problema de navegación concreto.

\subsection{Diseño de la recompensa}
Ng, Harada y Russell \cite{ng1999} probaron que el \emph{shaping} de forma potencial
$F=\gamma\,\Phi(s')-\Phi(s)$ no altera la política óptima. El resultado es correcto y
se cita como receta segura. Lo que reportamos es un caso práctico donde esa receta
falla: con $\gamma<1$ y distancias grandes, la garantía teórica sobrevive pero el
agente aprendido no, porque el residuo de descuento premia oscilar lejos de la meta
(Sec.~\ref{sec:disc}). No conocemos este efecto documentado con datos en un problema
de navegación sobre terreno real, y es uno de los aportes del trabajo.

\subsection{Políticas condicionadas a la meta}
Cuando la meta cambia entre episodios, el enfoque habitual es condicionar la función
de valor a la meta (Universal Value Function Approximators, UVFA)
\cite{schaul2015}, reetiquetar las trayectorias fallidas
(Hindsight Experience Replay, HER) \cite{andrychowicz2017} o priorizar el muestreo
del búfer (Prioritized Experience Replay, PER) \cite{schaul2016}. Nosotros no
aplicamos estas técnicas; medimos el problema que las motiva. Mostramos que con
metas aleatorias y DQN estándar los valores $Q$ divergen y la tasa de éxito no se
sostiene, lo que justifica nuestra decisión de alcance (meta fija) y define la línea
de trabajo futuro.

\subsection{Resumen del aporte}
Frente a lo anterior, este trabajo aporta tres cosas: un entorno de navegación sobre
terreno real con observación parcial e inicio aleatorio donde la tabla deja de ser
viable; una comparación controlada de las familias basada en valor y basada en
política contra el óptimo exacto de Dijkstra; y dos hallazgos negativos
documentados con datos, la trampa del \emph{shaping} potencial y la divergencia de
$Q$ con metas aleatorias.
